% TOP_PROBLEM: {"id": 5300074, "title": "Continuous extension of external-ray rotation number", "category_id": 11, "category": {"id": 11, "name": "dynamical_systems", "display_name": "Dynamical Systems", "description": "Problems about long-term behavior of deterministic systems, Hamiltonian dynamics, and periodic orbits.", "slug": "dynamical-systems", "order_index": 11, "created_at": "2026-07-31T15:26:25.671Z"}, "status": "partially_solved"}
% FIRST_POSTED: null
% ENTRY_KIND: statement_only
% Imported from: batch06.tex; UnsolvedMath ID 5300074
\documentclass[11pt]{article}
\usepackage[margin=1in]{geometry}
\usepackage{amsmath,amssymb,mathtools}
\usepackage{fontspec,unicode-math}
\setmainfont{Latin Modern Roman}
\setsansfont{Latin Modern Sans}
\setmathfont{Latin Modern Math}
\usepackage{xurl,hyperref}
\hypersetup{colorlinks=true,linkcolor=blue,urlcolor=blue}
\newcommand{\sref}[2]{\hyperlink{source:#1:#2}{[S#2]}}
\newcommand{\cataloguescope}[1]{\paragraph{Mathematical question.}#1}
\begin{document}
% UnsolvedMath ID 5300074; source catalogue code AMR-052-0074
\setcounter{section}{5300073}
\section{Continuous extension of external-ray rotation number}\label{5300074}
{\small\sffamily UnsolvedMath ID 5300074 · Dynamical Systems}
\subsection{Definitions and mathematical statement}

\paragraph{Shared notation.}$\mathbb N=\{1,2,\ldots\}$, and $\mathbb Z,\mathbb Q,\mathbb R,\mathbb C$ denote the integers, rationals, reals and complex numbers. Any other conventions are specified locally.


Fix $n\geq2$ and let $\mathcal P_n$ consist of monic polynomials
\[p(z)=z^n+a_{n-1}z^{n-1}+\cdots+a_1z,\qquad |a_1|\geq1,\]
with the coefficient topology. An external ray is a ray approaching infinity at an angle $t\in\mathbb R/\mathbb Z$ in the Böttcher coordinate, continued towards the Julia set wherever possible. On external angles $p$ acts by $t\mapsto nt$. When at least one external ray lands at $0$, the rays landing there have a well-defined circular rotation number. Connectedness of $J(p)$ is not assumed.
\paragraph{Statement recorded in the supplied source.}
Does this external-ray rotation number extend uniquely to a continuous map $\mathcal P_n\to\mathbb R/\mathbb Z$, taking the value $\theta$ whenever $a_1=e^{2\pi i\theta}$? \sref{5300074}{2}
\subsection{Short English statement}
Does the rotation number of rays landing at the fixed point extend uniquely and continuously over the full stated polynomial space, with its prescribed value on unit-modulus multipliers?
\subsection{Sources}
\begin{description}
\item[\hypertarget{source:5300074:1}{S1}] ULAM.AI and the UnsolvedMath Contributors, UnsolvedMath: A Curated Collection of Open Mathematics Problems, record 5300074 (AMR-052-0074). CC BY 4.0. \url{https://huggingface.co/datasets/ulamai/UnsolvedMath}.
\item[\hypertarget{source:5300074:2}{S2}] Ben Bielefeld (editor) and conference contributors, \emph{Conformal Dynamics Problem List}, arXiv:math/9201271v1 (1990), \S2, Question 12 and preceding paragraph, printed p. 8. \url{https://arxiv.org/abs/math/9201271v1}.
\end{description}
\label{end:5300074}
\end{document}
